\section{因子の選び方}\label{節:部分積分の因子選択}

前節の例題では,一次式 $x$ を微分し,指数関数 $\e^x$ を積分した.逆に選ぶと計算が重くなることも見た.では一般に,どちらの因子を微分し,どちらを積分すればよいのだろうか.目標は,残る積分 $\int u'v\dd x$ が既知の積分になるように選ぶことである.

まず,積分する側には条件がある.原始関数 $v$ を,その時点で求められなければならない.指数関数や三角関数は,積分しても同じ種類の関数のままなので,積分する側に選びやすい.一方,整式は微分すると次数が下がるため,微分する側の候補になる.

ただし,この順序は絶対の規則ではない.対数 $\log x$ の原始関数はまだ求めていないので,\,$\log x$ は積分する側には置けない.そこで $\log x$ を微分する側に選んでやる.\,$n\neq-1$ とし,
\[
 \int x^n\log x\dd x
\]
を考えよう.\,$u=\log x$,\,$v'=x^n$ と選ぶと,
\[
 u'=\frac1x,\qquad v=\frac{x^{n+1}}{n+1}
\]
である.部分積分の公式より,
\[
\begin{aligned}
 \int x^n\log x\dd x
 &=\frac{x^{n+1}}{n+1}\log x-\int\frac{x^{n+1}}{n+1}\cdot\frac1x\dd x\\
 &=\frac{x^{n+1}}{n+1}\log x-\frac{1}{n+1}\int x^n\dd x\\
 &=\frac{x^{n+1}}{n+1}\log x-\frac{x^{n+1}}{(n+1)^2}+C
\end{aligned}
\]
を得る.対数を微分して得た $\frac1x$ が,整式を積分して得た $x^{n+1}$ と約分され,残る積分が整式になったのだ.ここでは整式 $x^n$ の方を積分する側に選んでいる.関数の名前だけで順序を決めず,\,$u'v$ を具体的に見ることが肝である.

特に $n=0$ のときは $x^0=1$ なので,被積分関数は $\log x$ だけである.見た目が積でなくても,\,$\log x\cdot1$ と $1$ を掛けた積と考えれば,積分する側に既知の関数 $1$ を置けるのである.

\begin{bookpoint}
\textbf{POINT}\par
整式は微分する側,指数関数・三角関数は積分する側に選ぶ.対数は微分する側に選び,\,$\frac1x$ を残りの因子と約分する.
\end{bookpoint}

\begin{example}{1つの関数を積と見る}
$x>0$ において,\,$\int\log x\dd x$ を求めよ.
\end{example}

\paragraphbf{方針}
$\log x=\log x\cdot1$ と見て,\,$\log x$ を微分し,\,$1$ を積分する.

\paragraphbf{解答}
$x\log x-x+C$

\paragraphbf{解法}
$\log x$ の微分は $\frac1x$ であり,これに $1$ の原始関数 $x$ が掛かれば積は $1$ になる.そこで
\[
 u=\log x,\qquad v'=1,
 \qquad u'=\frac1x,\qquad v=x
\]
と選ぶと,
\[
\begin{aligned}
 \int\log x\dd x
 &=x\log x-\int x\cdot\frac1x\dd x\\
 &=x\log x-\int 1\dd x\\
 &=x\log x-x+C
\end{aligned}
\]
となる.答えを微分すると,
\[
 \frac{\dd}{\dd x}(x\log x-x)
 =\log x+x\cdot\frac1x-1
 =\log x
\]
である.

\paragraphbf{置換を先に考える場合}
上の一般形で除いた $n=-1$ の場合,すなわち $\int\frac{\log x}{x}\dd x$ では,\,$v=\frac{x^{n+1}}{n+1}$ が作れない.しかしこの場合は,前章で見たとおり $t=\log x$ と置くほうが短い.\,$\frac1x\dd x=\dd t$ となり,ただちに $\int t\dd t$ へ変わるからである.対数を見たら部分積分と決めるのではなく,対数の微分 $\frac1x$ がすでに残りの部分にあるかを先に調べよう.

積分する側の条件についても同じことがいえる.たとえば $\int x\e^{x^2}\dd x$ で,\,$x$ を微分しようとして $v'=\e^{x^2}$ と選ぶと,そこで原始関数が分からなくなる.しかし $t=x^2$ と置けば,\,$x\dd x=\frac12\dd t$ なので,
\[
 \int x\e^{x^2}\dd x
 =\frac12\int\e^t\dd t
 =\frac12\e^{x^2}+C
\]
である.部分積分を試す前に,置換に必要な微分の因子が含まれていないかを見ると,遠回りを避けられる.
